\begin{scanpage}{0011}{vii}
Therefore, it is natural to pose the question: in what way will the subsequent asymptotic equations differ from them when at last one obtains them in the mathematical theory of diffraction? One can say beforehand that the main term of the asymptotic expansion will not, in the general case, agree with the solution obtained on the basis of physical considerations: other (as a rule more complicated) slowly varying functions which determine the decay of the fields and currents as one goes further from the edges and borders, and also the diffraction interaction of the edges and the shadowing of the edge waves will figure in the main term. However, the refinement of the slowly varying functions in the expression for the diffraction field is not able to seriously influence the quantitative relationships. This is seen from a comparison of the results obtained in this book with calculations based on rigorous theory and other approximation equations, and also with the results of measurements.

The relationships obtained in this book also should help the development of asymptotic methods in the mathematical theory of diffraction, since they suggest the character of the approximations and the structure of the desired solution.
\vspace{1cm}\begin{flushright}L. A. Vaynshteyn\end{flushright}
\end{scanpage}
